\chapter{Polarity-Tuned Hydrogel Electrolytes}

\section{Introduction}

The liquid-electrolyte study identified polarity-related features as important
descriptors of solvent effectiveness, but it did not establish that bulk-liquid
relationships remain valid inside a hydrogel. Polymer confinement introduces
hydrogen bonding, local dielectric heterogeneity, restricted molecular motion,
preferential solvent uptake, and redox-ion transport through a dynamically
reorganizing network. These effects can alter both solvation entropy and ionic
mobility. The hydrogel study therefore asked whether a solvent polarity rule
derived from liquid electrolytes could be converted into a quantitative design
principle for a PVA matrix. Derivatives from the
dimethyl sulfoxide (DMSO), N,N-dimethylformamide (DMF), and
N-methyl-2-pyrrolidone (NMP) families were first compared within PVA
hydrogels. Data-driven analysis was then used to identify a cross-family
structure--property relationship. The predicted optimum was explored through
continuous polarity tuning, and the resulting hydrogel was evaluated from
molecular solvation to device-level output.

\section{Solvent-Family Library in PVA Hydrogels}

Within each solvent family, controlled substitution changed molecular
properties without replacing the complete chemical scaffold. Across families,
the persistence of a trend could be tested against different backbones.
Initial structural descriptors generated by cheminformatics tools were
sufficiently complex to fit the limited dataset but provided little
mechanistic guidance. Isolated HOMO and LUMO energies and electrostatic
potential maps likewise did not produce a consistent cross-family relationship
with thermopower.

Dipole moment provided a more integrated descriptor. Across the DMSO, DMF, and
NMP derivative families, $S_e$ first increased and then decreased with dipole
moment, yielding a reproducible volcano-shaped relationship. Weakly polar
solvents do not perturb the redox-ion solvation shells sufficiently. At the
opposite extreme, strongly polar solvents can over-stabilize and over-order the
local environment. Intermediate polarity provides a balance: the solvent
interacts strongly enough to reconstruct the two solvation shells, but not so
strongly that both become similarly constrained. The maximum thermopower
therefore occurs near an intermediate, rather than maximal, polarity
.

\section{Polarity-Tuned Hydrogel Electrolytes}

The discrete set of pure solvents could locate the optimum only coarsely.
Continuous tuning was achieved by combining methylsulfonylmethane (MSM) with a
lower-polarity co-solvent, denoted PD, in a PVA hydrogel. Varying the PD content
shifted the effective polarity through the region surrounding the predicted
volcano peak without changing the redox couple or polymer backbone.

Hydrogels containing 4, 6, 8, 10, and 12 vol\% PD were evaluated. The maximum
thermopower, 3.28 mV K$^{-1}$, occurred at 8 vol\% PD. Lower additions
produced values close to the MSM-containing hydrogel, whereas 12 vol\% PD
reduced performance. This composition dependence confirmed that the optimum
was not an artifact of comparing unrelated pure solvents and demonstrated why
a monotonic ``higher polarity is better'' rule would fail.

\section{Molecular Origin of the Enhanced Thermopower}

Cryogenic scanning electron microscopy, solution appearance, and X-ray
diffraction indicated that PD created local chemical heterogeneity in the
optimized hydrogel. Ab initio molecular dynamics simulations provided a
molecular interpretation. In the presence of MSM, PD formed short-range
interactions with $[\mathrm{Fe(CN)}_6]^{4-}$, whereas its contribution to the
first solvation shell of $[\mathrm{Fe(CN)}_6]^{3-}$ was negligible. The
coordination environment of the reduced ion was therefore reorganized more
strongly than that of the oxidized ion.

UV--visible and Raman spectroscopy independently supported this asymmetric
response. Spectral features associated with
$[\mathrm{Fe(CN)}_6]^{4-}$ shifted more strongly upon solvent modification,
whereas those of $[\mathrm{Fe(CN)}_6]^{3-}$ changed little. Together, the
simulations and experiments suggest that the optimized co-solvent composition
preferentially destabilized and reorganized the solvation shell of the more
highly charged redox ion. The increase in the difference between the two
temperature-dependent solvation free energies connects tuned polarity to
enhanced thermopower.

\section{Thermo-Electrochemical and Dynamic Performance}

The optimized composition exposed a trade-off that thermopower alone would
miss. Electrochemical impedance spectroscopy showed that the high-frequency
real-axis intercept increased from 322 $\Omega$ for the MSM hydrogel to
423 $\Omega$ for the optimized MSM--PD hydrogel, indicating a modest loss in
ionic transport. Nevertheless, the larger thermopower produced a net device
benefit. At $\Delta T=10$ K, the maximum power density increased from
0.24 to 0.44 $\mu$W cm$^{-2}$. The optimized thermocell delivered
0.11, 0.25, and 0.44 $\mu$W cm$^{-2}$ at temperature differences of
5, 7.5, and 10 K, respectively, while maintaining a normalized power density
of approximately 43 $\mu$W m$^{-2}$ K$^{-2}$.

Dynamic measurements were included because steady-state $S_e$ does not
describe how quickly a device follows a changing heat source. Under a constant
5 K gradient, the thermocell maintained an output of approximately 18 mV for
more than 1.5 h. It also responded reversibly during repeated temperature
cycling between 0 and 10 K. With a 5 k$\Omega$ load, quasi-continuous
charge--discharge cycling caused only a modest voltage decrease from
approximately 18 to 15 mV over 30 cycles, and the voltage recovered when the
temperature difference was stepped down and returned to 10 K. Thermal-response
kinetics and reversibility are therefore independent practical criteria, not
consequences that can be inferred from thermopower or ionic conductivity alone.

